\par\noindent\begin{minipage}{\linewidth}
% !TeX root = ../教材题库_全部条目_转写工作稿.tex
% Source: XJ-B1-C03-SU-P089-Y0854
% 第3章 / 小结与复习 小结与复习 / 教材 p.89 / PDF 第96页
% 已对照原页转写题干及公式；待教师复核；未制作答案
% 原题及相邻原解截图（核对用）：../题目原图/XJ-B1-C03-SU-P089-Y0854-01.webp
% 本文件只在转写完成后作为题干使用；OCR 条目尚不是可靠题干。
\begin{bankitem}{第3章 小结与复习 第1题}{第3章 / 小结与复习 小结与复习 / 教材 p.89 / PDF 第96页}{XJ-B1-C03-SU-P089-Y0854}{已对照原页转写题干及公式；待教师复核；未制作答案}
初中已学习函数的概念，在本章，我们进一步用更加严谨的集合和对应语言来刻画函数，突出了函数概念的本质是一个数集到另一个数集的一种确定的对应关系，明确了构成函数的要素，并引入函数符号 $y=f(x)$。通过本章的学习，你对函数概念有什么新的认识和体会吗？
\end{bankitem}
\end{minipage}\par\medskip
